\chapter{Testing}

Per la fase di testing sono state sviluppate quattro diverse implementazioni, una per ciascun tipo di dato precedentemente descritto, 
ciascuna organizzata all'interno della rispettiva cartella.

Sono stati condotti due test distinti: il primo analizza l’andamento delle prestazioni al crescere della dimensione dell’input, considerando 
valori pari alle potenze di 10; il secondo, invece, studia l’evoluzione degli stessi parametri al raddoppio della dimensione dell’input.

Nel seguito viene presentata la struttura e le scelte adottate per il test basato sulle potenze di 10; tuttavia, le considerazioni risultano 
valide anche per il test con incremento per raddoppio.

\section{Metriche}

Durante l'esecuzione dei test vengono generati due file distinti:

\begin{lstlisting}[language=C++, caption=File di output dei test]
std::pair<std::ofstream, std::ofstream> setup() {
    std::ofstream op_file("op_data.csv");
    std::ofstream hash_file("hash_data.csv");

    // Intestazione
    op_file << "OVERFLOW,HASH,RESIZE,Elements,Max_load_factor,Operazione,Latenza,Load_factor,Memoria,Key_distribution\n";
    hash_file << "OVERFLOW,HASH,RESIZE,Elements,Max_load_factor,Latenza,Memoria\n";

    return {std::move(op_file), std::move(hash_file)};
}
\end{lstlisting}

Il file \texttt{hash\_data.csv} contiene informazioni aggregate relative a ciascuna istanza di test, memorizzando le caratteristiche della configurazione 
utilizzata (strategia di gestione delle collisioni, funzione di hash, politica di resize, numero di elementi e fattore di carico massimo) 
insieme a metriche prestazionali globali quali latenza complessiva e memoria utilizzata.

Il file \texttt{op\_data.csv}, invece, raccoglie informazioni a grana fine per ogni singola operazione eseguita su ciascuna istanza. 
Per ogni operazione vengono registrati il tipo di istanza, la tipologia di operazione eseguita, la latenza, la memoria occupata, 
il fattore di carico corrente e un indice di distribuzione delle chiavi.

La latenza viene misurata come tempo CPU utilizzando la funzione \texttt{clock\_gettime} con clock \texttt{CLOCK\_PROCESS\_CPUTIME\_ID}, 
che consente di ottenere una misura stabile del tempo di esecuzione del processo, riducendo l'influenza di fattori esterni.

\subsection*{Distribuzione delle chiavi}

La distribuzione delle chiavi viene stimata mediante metriche differenti a seconda della strategia di gestione delle collisioni adottata:

\begin{itemize}
    \item \textbf{Open Addressing (Linear Probing e Double Hashing)}: viene calcolata la dimensione media dei cluster. Un cluster è definito 
    come una sequenza contigua di celle occupate. L'algoritmo identifica tali sequenze, ne calcola la lunghezza e restituisce la media tra tutte 
    le lunghezze osservate. Valori elevati indicano la presenza di fenomeni di clustering, generalmente associati a un peggioramento delle prestazioni.

    \item \textbf{Chaining}: viene calcolata la deviazione standard della lunghezza dei bucket. Questa misura consente di valutare quanto 
    uniformemente gli elementi siano distribuiti tra le liste: valori bassi indicano una distribuzione omogenea, mentre valori elevati 
    evidenziano squilibri tra i bucket.

    \item \textbf{Cuckoo Hashing}: viene adottato un approccio analogo all'open addressing, calcolando la dimensione media dei cluster su entrambe 
    le tabelle hash utilizzate dalla struttura. Le due tabelle vengono considerate congiuntamente nel processo di individuazione dei cluster.
\end{itemize}

Si osserva che le metriche adottate non sono direttamente confrontabili tra le diverse strategie, ma risultano utili per analizzare il comportamento 
interno di ciascun metodo.

\section{Struttura dei test}

Il test esegue una routine variando, per ogni iterazione, il valore del \texttt{load\_factor}. In particolare, vengono considerate tre configurazioni:
\begin{itemize}
    \item valore standard (consigliato per la struttura),
    \item valore inferiore,
    \item valore superiore.
\end{itemize}

\begin{lstlisting}[language=C++, caption=Test con differenti load factor]
std::pair<std::ofstream, std::ofstream> files = setup();

for(int x=0; x<5; x++) {
    // Standard load factor
    for(int i=1; i<5; i++) {
        std::cout << "-----------------------------" <<std::endl;
        std::cout << "ELEMENT COUNT: " << std::pow(10, i) <<std::endl;
        std::cout << "-----------------------------" <<std::endl;
        intTest(std::pow(10, i), files);
    }

    // Low load factor
    for(int i=1; i<5; i++) {
        std::cout << "-----------------------------" <<std::endl;
        std::cout << "ELEMENT COUNT: " << std::pow(10, i) <<std::endl;
        std::cout << "-----------------------------" <<std::endl;
        intTest(std::pow(10, i), files, 1);
    }

    // High load factor
    for(int i=1; i<5; i++) {
        std::cout << "-----------------------------" <<std::endl;
        std::cout << "ELEMENT COUNT: " << std::pow(10, i) <<std::endl;
        std::cout << "-----------------------------" <<std::endl;
        intTest(std::pow(10, i), files, 2);
    }
}
\end{lstlisting}

Ogni iterazione invoca una funzione di test che genera una nuova istanza della classe \texttt{HashTable} con differenti parametri di configurazione.

\begin{lstlisting}[language=C++, caption=Configurazione delle istanze]
/*
* HashTable:
*   OVERFLOW::OPEN_ADDRESSING_LINEAR_PROBING
*   HASH::MODE_1
*   RESIZE::DOUBLE
* 
* Key::INTEGER
* Value::INTEGER
*/
{
    std::cout << "-----------------------------" <<std::endl;
    std::cout << "HashTable:" << std::endl;
    std::cout << "  OVERFLOW::OPEN_ADDRESSING_LINEAR_PROBING" << std::endl;
    std::cout << "  HASH::MODE_1" << std::endl;
    std::cout << "  RESIZE::DOUBLE" << std::endl;
    std::cout << "  Key::INTEGER" << std::endl;
    std::cout << "  Value::INTEGER" << std::endl;
    std::cout << "-----------------------------" <<std::endl;

    HashTable<int, int> HashTable(OVERFLOW::OPEN_ADDRESSING_LINEAR_PROBING, HASH::MODE_1, RESIZE::DOUBLE, 256, linear_load_factor);
    intTestRoutine(HashTable, elements, files.first, files.second);
}
{
    std::cout << "-----------------------------" <<std::endl;
    std::cout << "HashTable:" << std::endl;
    std::cout << "  OVERFLOW::OPEN_ADDRESSING_LINEAR_PROBING" << std::endl;
    std::cout << "  HASH::MODE_1" << std::endl;
    std::cout << "  RESIZE::DOUBLE" << std::endl;
    std::cout << "  Key::INTEGER" << std::endl;
    std::cout << "  Value::INTEGER" << std::endl;
    std::cout << "-----------------------------" <<std::endl;

    HashTable<int, int> HashTable(OVERFLOW::OPEN_ADDRESSING_LINEAR_PROBING, HASH::MODE_1, RESIZE::DOUBLE, 256, linear_load_factor);
    intTestRoutine(HashTable, elements, files.first, files.second, false);
}
\end{lstlisting}

La routine di test procede inizialmente alla generazione dei dati da inserire e successivamente esegue una sequenza di operazioni selezionate 
in modo pseudo-casuale secondo la seguente distribuzione:

\begin{itemize}
    \item 40\% \texttt{Exist}
    \item 30\% \texttt{Insert}
    \item 10\% \texttt{Update}
    \item 10\% \texttt{Get}
    \item 10\% \texttt{Remove}
\end{itemize}

Durante l'esecuzione, la routine può operare in due modalità:
\begin{itemize}
    \item \textbf{Modalità analitica}: ogni operazione viene registrata nel file \texttt{op\_data.csv};
    \item \textbf{Modalità prestazionale}: le operazioni vengono eseguite senza logging dettagliato, al fine di evitare overhead e 
    misurare accuratamente le prestazioni complessive, salvate in \texttt{hash\_data.csv}.
\end{itemize}

\begin{lstlisting}[language=C++, caption=Test routine]
template <typename K, typename V>
void intTestRoutine(HashTable<K, V>& HashTable, int elements, std::ofstream& file_hash, std::ofstream& file_op, bool mode = true) {
    const int ELEMENTS = elements;
    int errors = 0;
    // Generatore di numeri casuali
    std::vector<std::pair<int, int>> numbers;
    numbers.reserve(ELEMENTS);

    std::random_device rd;
    std::mt19937 gen(rd());
    std::uniform_int_distribution<> dist(1, ELEMENTS*10);

    int value = 0;
    for (int i = 0; i < ELEMENTS; ++i) {
        value = dist(gen);
        numbers.emplace_back(value, value);
    }

    struct timespec start, end;
    if(!mode) {
        clock_gettime(CLOCK_PROCESS_CPUTIME_ID, &start);
    }

    std::size_t counter = 0;
    while(counter < numbers.size()) {
        std::uniform_int_distribution<> dist2(0, 9);
        switch(dist2(gen)) {
            case 0:
            case 1:
            case 2:
            case 3: {
                // Exist
                if(counter > 0) {
                    std::uniform_int_distribution<> dist_temp(0, numbers.size() - 1);
                    exist(HashTable, numbers[dist_temp(gen)].first, file_op, elements, mode);
                }
                break;
            }
            case 4:
            case 5:
            case 6: {
                // Insert
                insert(HashTable, numbers[counter], file_op, elements, mode);
                counter++;
                break;
            }
            case 7: {
                // Update
                if(counter > 0) {
                    std::uniform_int_distribution<> dist_temp(0, numbers.size() - 1);
                    update(HashTable, std::make_pair(numbers[dist_temp(gen)].first, dist(gen)), file_op, elements, mode);
                }
                break;
            }
            case 8: {
                // Get
                if(counter > 0) {
                    std::uniform_int_distribution<> dist_temp(0, numbers.size() - 1);
                    get(HashTable, numbers[dist_temp(gen)].first, file_op, elements, mode);
                }
                break;
            }
            case 9: {
                // Remove
                if(counter > 0) {
                    std::uniform_int_distribution<> dist_temp(0, numbers.size() - 1);
                    remove(HashTable, numbers[dist_temp(gen)].first, file_op, elements, mode);
                }
                break;
            }
        }
    }

    if(!mode) {
        clock_gettime(CLOCK_PROCESS_CPUTIME_ID, &end);
        double cpu_time = (end.tv_sec - start.tv_sec) + (end.tv_nsec - start.tv_nsec) / 1e9;

        csvHashPrint(
            file_hash,
            HashTable.getOverflow(),
            HashTable.getHash(),
            HashTable.getResize(),
            elements,
            HashTable.getMaxLoadFactor(),
            cpu_time,
            HashTable.maxMemoryUsage()
        );
    }
}
\end{lstlisting}

Ogni operazione eseguita sulla struttura (inserimento, ricerca, aggiornamento e rimozione) misura il tempo di esecuzione e, se abilitato, 
registra i risultati nel file \texttt{op\_data.csv}, includendo anche metriche dinamiche come il fattore di carico e la distribuzione delle chiavi.

\begin{lstlisting}[language=C++, caption=Esempio operazione]
template <typename K, typename V>
void insert(HashTable<K, V>& hashTable, const std::pair<K, V>& arr, std::ofstream& file, int elements, bool print = true) {
    if(print) {
        struct timespec start, end;
        clock_gettime(CLOCK_PROCESS_CPUTIME_ID, &start);

        hashTable.insert(arr.first, arr.second);

        clock_gettime(CLOCK_PROCESS_CPUTIME_ID, &end);
        double cpu_time = (end.tv_sec - start.tv_sec) + (end.tv_nsec - start.tv_nsec) / 1e9;

        csvOperationPrint(
            file,
            hashTable.getOverflow(),
            hashTable.getHash(),
            hashTable.getResize(),
            elements,
            hashTable.getMaxLoadFactor(),
            "INSERT",
            cpu_time,
            hashTable.getLoadFactor(),
            hashTable.memoryUsage(),
            hashTable.getKeyDistribution()
        );
    } else {
        hashTable.insert(arr.first, arr.second);
        hashTable.memoryUsage();
    }
}
\end{lstlisting}

Il build dell'apparato di testing viene seguito con i seguenti parametri di ottimizzazione

\begin{lstlisting}[language=bash, caption=Parametri compilazione]
g++ -O3 -march=native -flto -std=c++20 intTest.cpp -o ./build/intTest 
\end{lstlisting}